\chapter{Discussion and Conclusions}

\section{Synthesis of Results and Comparative Analysis}
This investigation successfully established a transient performance baseline 
for sustainable polymer alternatives under the EN 397 regulatory framework. 
The transmitted force results (2.25--4.54 kN) indicate that all 
candidate materials meet the primary safety criterion of < 5.0 kN. 

The filtered Virgin HDPE baseline of 4.05 kN is consistent with 
experimental ranges of 3.0--4.5 kN reported by Tham et al. \cite{tham2014} 
and safety researchers for standard Type I shells. This alignment 
suggests that the modeling of the shell geometry and contact 
interactions is physically representative of the standard shock 
absorption test. The slight elevation in peak force compared to 
benchmark FEA studies by Long et al. \cite{long2015} (3.2--3.8 kN) 
may be attributed to the rigid headform assumption, which neglects the 
compliance of the EN 960 headform assembly, thereby providing a 
conservative assessment of force transmission.

\section{Energy Absorption and Protective Mechanisms}
The results indicate two distinct mechanical strategies for impact 
protection. The highly ductile Virgin HDPE neutralises impact energy 
primarily through internal strain energy yielding (27.35 J). This 
mechanism effectively extends the impact duration, thereby attenuating 
the peak transmitted force. However, this ductile response is 
associated with the highest structural crush depth (18.42 mm), which 
limits the available suspension travel in high-velocity scenarios.

Conversely, the high-modulus materials (ABS, rPC) protect the wearer 
via load distribution due to increased flexural stiffness. By 
distributing the impact force across a broader surface area of the 
EPS liner, these materials transmit lower localised stress to the 
shell structure. The superior performance of the ABS shell (2.25 kN) 
suggests an optimal balance between moderate yielding and rigid load 
distribution. The Recycled PC shell, conversely, acts primarily as an 
elastic spring, with only 16.67 J absorbed as internal strain energy. 
The high elastic rebound (54.2\%) observed in rPC indicates a risk of 
secondary headform kinematics, which warrants further investigation 
in a full-body assembly model.

\section{Rotational Kinematics and Injury Risk Correlation}
The 30$^\circ$ oblique impact simulations (Phase 5) suggest that 
material modulus is a primary factor in rotational acceleration 
generation. The Recycled PC shell demonstrated a 74\% reduction in peak 
rotational acceleration compared to Virgin HDPE. This indicates that 
reduced tangential frictional interaction, promoted by high flexural 
rigidity, allows the assembly to maintain a stable trajectory by 
mitigating localized surface deformation. 

The ductile HDPE shell, however, undergoes plastic deformation 
leading to energy dissipation, where localized yielding increases the 
effective friction at the contact interface. This results in higher 
rotational torque, which is associated with an increased risk of 
Diffuse Axonal Injury (DAI) \cite{mills2011}. The extreme rotational 
acceleration recorded for the Bamboo-HDPE composite (12044 rad/s$^2$) 
suggests that the presence of stiff fibres in a ductile matrix may 
artificially increase surface 'grab' without providing the bulk 
rigidity required to promote a low-friction response. 

\section{Discussion of Limitations and Validation Confidence}
\label{sec:limitations}

The following limitations are inherent to the modelling assumptions. The 
overall validation confidence is considered high for relative material 
ranking, though absolute magnitudes are subject to the following 
uncertainties:
\begin{itemize}
    \item \textbf{Isotropic assumption for Bamboo-HDPE.} The 
    isotropic modelling of the 30\% wt composite likely overestimates 
    stiffness in transverse directions. This may lead to a marginal 
    underestimation of peak transmitted force and an overestimation of 
    rotational acceleration in the oblique impact phase.
    \item \textbf{EPS liner simplification.} Modeling the EPS as linear 
    elastic likely overestimates the transmitted force by 10--15\%, as the 
    energy-dissipating cell crushing plateau is neglected; the 
    EN 397 passes reported here likely represent a conservative safety 
    margin.
    \item \textbf{Rigid headform constraint.} The exclusion of assembly 
    compliance marginally increases recorded force peaks. However, this 
    ensures that the comparison is focused strictly on the shell material 
    performance rather than assembly-specific dynamics.
    \item \textbf{Numerical signal processing.} While the SAE J211 CFC 60 
    filter is required to resolve meaningful force histories from explicit 
    penalty contact, the 100 Hz cutoff may smooth transient peaks that occur 
    at the millisecond scale. 
\end{itemize}

\section{Conclusions}
Sustainable polymer alternatives, particularly Recycled PC and 
Bamboo-HDPE 30\% wt, demonstrate structural compliance with the EN 397 
force transmission standard. While Recycled PC offers superior 
protection against rotational acceleration in oblique impacts, 
Bamboo-HDPE composites may require surface coating or geometric 
modifications to mitigate rotational injury risk. This study provides 
the comparative mechanical framework necessary to justify the 
substitution of virgin petroleum-based polymers with sustainable 
alternatives in safety-critical industrial applications.

\section{Future Work}
Future research should focus on the following targeted, actionable gaps:
\begin{enumerate}
    \item Perform in-house physical EN 397 drop tests to provide direct 
    validation of the FEA material models.
    \item Extend simulation to optional EN 397 tests (e.g., 15$^\circ$ lateral 
    impacts).
    \item Model Bamboo-HDPE as an orthotropic material with 
    direction-dependent properties to assess the error introduced by the 
    isotropic assumption.
    \item Conduct a parametric study of wall thickness (2.5--4.0 mm) for the 
    sustainable materials to find the minimum compliant thickness.
    \item Perform an indicative LCA using Ecoinvent data to quantify the 
    CO$_2$e savings of sustainable material substitution.
\end{enumerate}

